\documentclass[utf8, 12pt]{beamer}

% === КОДИРОВКИ И ЯЗЫК ===
\usepackage[utf8]{inputenc}
\usepackage[T1,T2A]{fontenc}
\usepackage[russian]{babel}
\usepackage{lmodern}

% === ГРАФИКА ===
\usepackage{graphicx}
\usepackage{tikz}
\usetikzlibrary{shapes, arrows, positioning}

% === ТАБЛИЦЫ ===
\usepackage{booktabs}
\usepackage{array}

% === МАТЕМАТИКА ===
\usepackage{amsmath}
\usepackage{amssymb}

% === ЦВЕТА ===
\usepackage{xcolor}

% === ТЕКСТ ===
\usepackage{lipsum}


% === ЛИСТИНГИ КОДА ===
\usepackage{listings}
\lstset{
    basicstyle=\ttfamily\footnotesize,
    breaklines=true,
    frame=single,
    language=[LaTeX]TeX,
    showstringspaces=false,
    commentstyle=\color{green!50!black},
    keywordstyle=\color{blue},
    stringstyle=\color{red},
    extendedchars=true,
    inputencoding=utf8,
    literate={а}{{\selectfont\char224}}1
             {б}{{\selectfont\char225}}1
             {в}{{\selectfont\char226}}1
             {г}{{\selectfont\char227}}1
             {д}{{\selectfont\char228}}1
             {е}{{\selectfont\char229}}1
             {ё}{{\"e}}1
             {ж}{{\selectfont\char230}}1
             {з}{{\selectfont\char231}}1
             {и}{{\selectfont\char232}}1
             {й}{{\selectfont\char233}}1
             {к}{{\selectfont\char234}}1
             {л}{{\selectfont\char235}}1
             {м}{{\selectfont\char236}}1
             {н}{{\selectfont\char237}}1
             {о}{{\selectfont\char238}}1
             {п}{{\selectfont\char239}}1
             {р}{{\selectfont\char240}}1
             {с}{{\selectfont\char241}}1
             {т}{{\selectfont\char242}}1
             {у}{{\selectfont\char243}}1
             {ф}{{\selectfont\char244}}1
             {х}{{\selectfont\char245}}1
             {ц}{{\selectfont\char246}}1
             {ч}{{\selectfont\char247}}1
             {ш}{{\selectfont\char248}}1
             {щ}{{\selectfont\char249}}1
             {ъ}{{\selectfont\char250}}1
             {ы}{{\selectfont\char251}}1
             {ь}{{\selectfont\char252}}1
             {э}{{\selectfont\char253}}1
             {ю}{{\selectfont\char254}}1
             {я}{{\selectfont\char255}}1
             {А}{{\selectfont\char192}}1
             {Б}{{\selectfont\char193}}1
             {В}{{\selectfont\char194}}1
             {Г}{{\selectfont\char195}}1
             {Д}{{\selectfont\char196}}1
             {Е}{{\selectfont\char197}}1
             {Ё}{{\"E}}1
             {Ж}{{\selectfont\char198}}1
             {З}{{\selectfont\char199}}1
             {И}{{\selectfont\char200}}1
             {Й}{{\selectfont\char201}}1
             {К}{{\selectfont\char202}}1
             {Л}{{\selectfont\char203}}1
             {М}{{\selectfont\char204}}1
             {Н}{{\selectfont\char205}}1
             {О}{{\selectfont\char206}}1
             {П}{{\selectfont\char207}}1
             {Р}{{\selectfont\char208}}1
             {С}{{\selectfont\char209}}1
             {Т}{{\selectfont\char210}}1
             {У}{{\selectfont\char211}}1
             {Ф}{{\selectfont\char212}}1
             {Х}{{\selectfont\char213}}1
             {Ц}{{\selectfont\char214}}1
             {Ч}{{\selectfont\char215}}1
             {Ш}{{\selectfont\char216}}1
             {Щ}{{\selectfont\char217}}1
             {Ъ}{{\selectfont\char218}}1
             {Ы}{{\selectfont\char219}}1
             {Ь}{{\selectfont\char220}}1
             {Э}{{\selectfont\char221}}1
             {Ю}{{\selectfont\char222}}1
             {Я}{{\selectfont\char223}}1
}

% === ТЕМА BEAMER ===
\usetheme{Madrid}
\usecolortheme{default}

% === НАСТРОЙКИ ===
\setbeamertemplate{navigation symbols}{}
\setbeamertemplate{footline}[frame number]

% === ИНФОРМАЦИЯ О ПРЕЗЕНТАЦИИ ===
\title{Работа с изображениями и ссылками в \LaTeX}
\subtitle{Лабораторная работа №4}
\author{Викторов Егор}
\institute{Российский Университет Дружбы Народов имени Патриса Лумумбы}
\date{\today}

\begin{document}

% ==================== ТИТУЛЬНЫЙ СЛАЙД ====================
\begin{frame}
    \titlepage
\end{frame}

% ==================== СОДЕРЖАНИЕ ====================
\begin{frame}{Содержание}
    \tableofcontents
\end{frame}

% ==================== РАЗДЕЛ 1: ВВЕДЕНИЕ ====================
\section{Введение}

\begin{frame}{Цели работы}
    \begin{block}{Основные задачи}
        \begin{itemize}
            \item Изучить параметры команды \texttt{\textbackslash includegraphics}
            \item Исследовать поведение изображений с \texttt{textwidth} и \texttt{linewidth}
            \item Изучить работу плавающих объектов (floats)
            \item Понять механизм работы меток и ссылок (\texttt{label/ref})
            \item Экспериментировать с нумерацией разделов
        \end{itemize}
    \end{block}
\end{frame}

% ==================== РАЗДЕЛ 2: ПАРАМЕТРЫ ИЗОБРАЖЕНИЙ ====================
\section{Параметры изображений}

\begin{frame}{Основные параметры includegraphics}
    \begin{table}
        \centering
        \begin{tabular}{@{}lp{6cm}@{}}
            \toprule
            \textbf{Параметр} & \textbf{Описание} \\
            \midrule
            \texttt{width} & Ширина изображения \\
            \texttt{height} & Высота изображения \\
            \texttt{scale} & Масштаб (множитель) \\
            \texttt{angle} & Угол поворота (градусы) \\
            \texttt{keepaspectratio} & Сохранение пропорций \\
            \bottomrule
        \end{tabular}
    \end{table}
    
    \vspace{0.5cm}
    
    \begin{exampleblock}{Синтаксис}
        \texttt{\textbackslash includegraphics[параметры]\{файл\}}
    \end{exampleblock}
\end{frame}

% ==================== ДЕМОНСТРАЦИЯ WIDTH ====================
\begin{frame}{Параметр width: визуализация}
    \begin{center}
        \begin{tikzpicture}[scale=0.8]
            % 30% ширины
            \draw[fill=blue!20, draw=black, thick] (0,0) rectangle (3,2);
            \node at (1.5,1) {\textbf{30\%}};
            \node[below] at (1.5,-0.3) {\scriptsize\texttt{width=0.3\textbackslash textwidth}};
            
            % 50% ширины
            \draw[fill=green!20, draw=black, thick] (4,0) rectangle (9,2);
            \node at (6.5,1) {\textbf{50\%}};
            \node[below] at (6.5,-0.3) {\scriptsize\texttt{width=0.5\textbackslash textwidth}};
            
            % 80% ширины
            \draw[fill=orange!20, draw=black, thick] (0,-3) rectangle (8,-1);
            \node at (4,-2) {\textbf{80\%}};
            \node[below] at (4,-3.3) {\scriptsize\texttt{width=0.8\textbackslash textwidth}};
        \end{tikzpicture}
    \end{center}
    
    \begin{alertblock}{Рекомендация}
        Используйте относительные размеры для адаптивности!
    \end{alertblock}
\end{frame}

% ==================== ДЕМОНСТРАЦИЯ HEIGHT ====================
\begin{frame}{Параметр height: визуализация}
    \begin{center}
        \begin{tikzpicture}[scale=0.7]
            % Разные высоты
            \draw[fill=red!20, draw=black, thick] (0,0) rectangle (2,1);
            \node at (1,0.5) {\small 3cm};
            \node[below] at (1,-0.3) {\tiny\texttt{height=3cm}};
            
            \draw[fill=yellow!40, draw=black, thick] (3,0) rectangle (5,2);
            \node at (4,1) {\small 5cm};
            \node[below] at (4,-0.3) {\tiny\texttt{height=5cm}};
            
            \draw[fill=cyan!30, draw=black, thick] (6,0) rectangle (8,3);
            \node at (7,1.5) {\small 8cm};
            \node[below] at (7,-0.3) {\tiny\texttt{height=8cm}};
            
            \draw[fill=purple!20, draw=black, thick] (9,0) rectangle (11,4);
            \node at (10,2) {\small 10cm};
            \node[below] at (10,-0.3) {\tiny\texttt{height=10cm}};
        \end{tikzpicture}
    \end{center}
\end{frame}

% ==================== ДЕМОНСТРАЦИЯ SCALE ====================
\begin{frame}{Параметр scale: визуализация}
    \begin{center}
        \begin{tikzpicture}[scale=0.6]
            % Базовый размер
            \draw[fill=gray!30, draw=black, thick] (0,0) rectangle (2,1.5);
            \node at (1,0.75) {1.0};
            \node[below] at (1,-0.5) {\scriptsize\texttt{scale=1.0}};
            
            % Увеличенный
            \draw[fill=gray!30, draw=black, thick] (3.5,0) rectangle (6.5,2.25);
            \node at (5,1.125) {1.5};
            \node[below] at (5,-0.5) {\scriptsize\texttt{scale=1.5}};
            
            % Ещё больше
            \draw[fill=gray!30, draw=black, thick] (8,0) rectangle (12,3);
            \node at (10,1.5) {2.0};
            \node[below] at (10,-0.5) {\scriptsize\texttt{scale=2.0}};
            
            % Уменьшенный
            \draw[fill=gray!30, draw=black, thick] (0,-3) rectangle (1,-2.25);
            \node at (0.5,-2.625) {\tiny 0.5};
            \node[below] at (0.5,-3.3) {\scriptsize\texttt{scale=0.5}};
        \end{tikzpicture}
    \end{center}
    
    \begin{block}{Примечание}
        \texttt{scale} умножает оригинальный размер изображения
    \end{block}
\end{frame}

% ==================== ДЕМОНСТРАЦИЯ ANGLE ====================
\begin{frame}{Параметр angle: визуализация}
    \begin{center}
        \begin{tikzpicture}[scale=0.7]
            % 0 градусов
            \begin{scope}
                \draw[fill=teal!30, draw=black, thick] (0,0) rectangle (2,1.5);
                \node at (1,0.75) {$0°$};
                \node[below] at (1,-0.5) {\scriptsize\texttt{angle=0}};
            \end{scope}
            
            % 45 градусов
            \begin{scope}[xshift=4cm]
                \begin{scope}[rotate around={45:(1,0.75)}]
                    \draw[fill=teal!30, draw=black, thick] (0,0) rectangle (2,1.5);
                    \node at (1,0.75) {$45°$};
                \end{scope}
                \node[below] at (1,-1) {\scriptsize\texttt{angle=45}};
            \end{scope}
            
            % 90 градусов
            \begin{scope}[xshift=8cm]
                \begin{scope}[rotate around={90:(1,0.75)}]
                    \draw[fill=teal!30, draw=black, thick] (0,0) rectangle (2,1.5);
                    \node at (1,0.75) {$90°$};
                \end{scope}
                \node[below] at (1,-0.5) {\scriptsize\texttt{angle=90}};
            \end{scope}
            
            % 180 градусов
            \begin{scope}[xshift=12cm]
                \begin{scope}[rotate around={180:(1,0.75)}]
                    \draw[fill=teal!30, draw=black, thick] (0,0) rectangle (2,1.5);
                    \node at (1,0.75) {$180°$};
                \end{scope}
                \node[below] at (1,-0.5) {\scriptsize\texttt{angle=180}};
            \end{scope}
        \end{tikzpicture}
    \end{center}
    
    \begin{alertblock}{Важно!}
        Поворот происходит \textbf{против часовой стрелки}
    \end{alertblock}
\end{frame}

% ==================== КОМБИНАЦИЯ ПАРАМЕТРОВ ====================
\begin{frame}{Комбинирование параметров}
    \begin{center}
        \begin{tikzpicture}[scale=0.8]
            % Оригинал
            \draw[fill=lime!30, draw=black, thick] (0,0) rectangle (2,1.5);
            \node at (1,0.75) {\small Оригинал};
            \node[below] at (1,-0.5) {\tiny без параметров};
            
            % Width + angle
            \begin{scope}[xshift=4cm]
                \begin{scope}[rotate around={30:(1.5,1.125)}]
                    \draw[fill=lime!30, draw=black, thick] (0,0) rectangle (3,2.25);
                    \node at (1.5,1.125) {\small Width+Angle};
                \end{scope}
                \node[below] at (1.5,-1) {\tiny\texttt{width=..., angle=30}};
            \end{scope}
            
            % Scale + angle
            \begin{scope}[xshift=9cm]
                \begin{scope}[rotate around={-20:(1,0.75)}, scale=1.5]
                    \draw[fill=lime!30, draw=black, thick] (0,0) rectangle (2,1.5);
                    \node at (1,0.75) {\tiny Scale+Angle};
                \end{scope}
                \node[below] at (1.5,-1) {\tiny\texttt{scale=1.5, angle=-20}};
            \end{scope}
        \end{tikzpicture}
    \end{center}
    
    \begin{exampleblock}{Пример кода}
        \texttt{\textbackslash includegraphics[width=0.5\textbackslash textwidth, angle=30]\{image.png\}}
    \end{exampleblock}
\end{frame}

% ==================== РАЗДЕЛ 3: TEXTWIDTH VS LINEWIDTH ====================
\section{textwidth vs linewidth}

\begin{frame}{Разница между textwidth и linewidth}
    \begin{columns}[T]
        \begin{column}{0.48\textwidth}
            \begin{block}{\texttt{\textbackslash textwidth}}
                \begin{itemize}
                    \item Полная ширина текстовой области
                    \item Не меняется внутри окружений
                    \item Постоянная величина
                \end{itemize}
            \end{block}
        \end{column}
        
        \begin{column}{0.48\textwidth}
            \begin{block}{\texttt{\textbackslash linewidth}}
                \begin{itemize}
                    \item Текущая ширина строки
                    \item Меняется в колонках, списках
                    \item Адаптивная величина
                \end{itemize}
            \end{block}
        \end{column}
    \end{columns}
    
    \vspace{0.5cm}
    
    \begin{alertblock}{Важно!}
        В режиме \texttt{twocolumn} значения различаются!
    \end{alertblock}
\end{frame}

% ==================== ВИЗУАЛИЗАЦИЯ TEXTWIDTH VS LINEWIDTH ====================
\begin{frame}{Визуализация: textwidth vs linewidth}
    \begin{center}
        \begin{tikzpicture}[scale=0.9]
            % Одна колонка
            \draw[thick] (0,0) rectangle (10,1.5);
            \node at (5,1.8) {\textbf{Одна колонка}};
            \draw[<->, red, thick] (0,-0.3) -- (10,-0.3);
            \node[red] at (5,-0.6) {\texttt{\textbackslash textwidth = \textbackslash linewidth}};
            
            % Две колонки
            \begin{scope}[yshift=-3cm]
                \draw[thick] (0,0) rectangle (4.5,1.5);
                \draw[thick] (5.5,0) rectangle (10,1.5);
                \node at (5,1.8) {\textbf{Две колонки (twocolumn)}};
                
                \draw[<->, red, thick] (0,-0.3) -- (10,-0.3);
                \node[red] at (5,-0.6) {\texttt{\textbackslash textwidth}};
                
                \draw[<->, blue, thick] (0,-0.9) -- (4.5,-0.9);
                \node[blue] at (2.25,-1.2) {\texttt{\textbackslash linewidth}};
                
                \draw[<->, blue, thick] (5.5,-0.9) -- (10,-0.9);
                \node[blue] at (7.75,-1.2) {\texttt{\textbackslash linewidth}};
            \end{scope}
        \end{tikzpicture}
    \end{center}
\end{frame}

% ==================== ПРИМЕР КОДА БЕЗ РУССКИХ КОММЕНТАРИЕВ ====================
\begin{frame}[fragile]{Пример использования}
    \begin{exampleblock}{Одна колонка}
\begin{lstlisting}
\documentclass{article}
\begin{document}
\includegraphics[width=0.8\textwidth]{image.png}
\includegraphics[width=0.8\linewidth]{image.png}
\end{document}
\end{lstlisting}
    \end{exampleblock}
    
    \vspace{0.3cm}
    \textbf{Примечание:} В одноколоночном режиме \texttt{textwidth} и \texttt{linewidth} дают одинаковый результат!
\end{frame}

\begin{frame}[fragile]{Пример с twocolumn}
    \begin{exampleblock}{Две колонки}
\begin{lstlisting}
\documentclass[twocolumn]{article}
\begin{document}
\includegraphics[width=0.8\textwidth]{img.png}

\includegraphics[width=0.8\linewidth]{img.png}
\end{document}
\end{lstlisting}
    \end{exampleblock}
    
    \vspace{0.3cm}
    \begin{alertblock}{Важно!}
        С \texttt{textwidth} изображение выйдет за пределы колонки! \\
        С \texttt{linewidth} --- правильно впишется в колонку.
    \end{alertblock}
\end{frame}

% ==================== РАЗДЕЛ 4: ПЛАВАЮЩИЕ ОБЪЕКТЫ ====================
\section{Плавающие объекты (floats)}

\begin{frame}{Что такое плавающие объекты?}
    \begin{block}{Определение}
        Плавающие объекты (floats) --- это элементы, которые \LaTeX\ может перемещать для оптимального размещения на странице.
    \end{block}
    
    \vspace{0.5cm}
    
    \begin{columns}[T]
        \begin{column}{0.48\textwidth}
            \textbf{Типы floats:}
            \begin{itemize}
                \item \texttt{figure} --- рисунки
                \item \texttt{table} --- таблицы
            \end{itemize}
        \end{column}
        
        \begin{column}{0.48\textwidth}
            \textbf{Преимущества:}
            \begin{itemize}
                \item Автоматическая нумерация
                \item Список рисунков/таблиц
                \item Оптимальное размещение
            \end{itemize}
        \end{column}
    \end{columns}
\end{frame}

% ==================== СПЕЦИФИКАТОРЫ ПОЗИЦИИ ====================
\begin{frame}{Спецификаторы позиции}
    \begin{table}
        \centering
        \begin{tabular}{@{}clp{5cm}@{}}
            \toprule
            \textbf{Спец.} & \textbf{Название} & \textbf{Описание} \\
            \midrule
            \texttt{h} & here & Примерно здесь \\
            \texttt{t} & top & Вверху страницы \\
            \texttt{b} & bottom & Внизу страницы \\
            \texttt{p} & page & На отдельной странице \\
            \texttt{!} & override & Игнорировать ограничения \\
            \texttt{H} & HERE & Строго здесь (пакет \texttt{float}) \\
            \bottomrule
        \end{tabular}
    \end{table}
    
    \vspace{0.3cm}
    
    \begin{exampleblock}{Примеры}
        \texttt{[htbp]} --- попробовать h, затем t, затем b, затем p \\
        \texttt{[!h]} --- здесь, игнорируя ограничения \\
        \texttt{[H]} --- строго здесь (требует \texttt{\textbackslash usepackage\{float\}})
    \end{exampleblock}
\end{frame}

% ==================== ВЗАИМОДЕЙСТВИЕ СПЕЦИФИКАТОРОВ ====================
\begin{frame}{Взаимодействие спецификаторов}
    \begin{center}
        \begin{tikzpicture}[
            node distance=1.5cm,
            box/.style={rectangle, draw, fill=blue!20, text width=2.5cm, align=center, minimum height=1cm}
        ]
            \node[box] (h) {\texttt{[h]}};
            \node[box, right=of h] (t) {\texttt{[t]}};
            \node[box, right=of t] (b) {\texttt{[b]}};
            \node[box, right=of b] (p) {\texttt{[p]}};
            
            \draw[->, thick] (h) -- (t) node[midway, above] {если не подходит};
            \draw[->, thick] (t) -- (b) node[midway, above] {если не подходит};
            \draw[->, thick] (b) -- (p) node[midway, above] {если не подходит};
        \end{tikzpicture}
    \end{center}
    
    \vspace{0.5cm}
    
    \begin{alertblock}{Важно!}
        \LaTeX\ пробует варианты \textbf{слева направо} в порядке указания
    \end{alertblock}
    
    \begin{block}{Рекомендация}
        Используйте \texttt{[htbp]} для максимальной гибкости
    \end{block}
\end{frame}

% ==================== ПРИМЕР FLOAT ====================
\begin{frame}[fragile]{Пример плавающего объекта}
    \begin{exampleblock}{Код}
\begin{lstlisting}
\begin{figure}[htbp]
    \centering
    \includegraphics[width=0.6\textwidth]{image.png}
    \caption{Podpis k risunku}
    \label{fig:example}
\end{figure}

Ssylka na risunok: sm. ris.~\ref{fig:example}.
\end{lstlisting}
    \end{exampleblock}
    
    \begin{alertblock}{Порядок важен!}
        \texttt{\textbackslash label} должна идти \textbf{после} \texttt{\textbackslash caption}!
    \end{alertblock}
\end{frame}

% ==================== РАЗДЕЛ 5: МЕТКИ И ССЫЛКИ ====================
\section{Метки и ссылки}

\begin{frame}{Система меток и ссылок}
    \begin{block}{Основные команды}
        \begin{itemize}
            \item \texttt{\textbackslash label\{метка\}} --- создать метку
            \item \texttt{\textbackslash ref\{метка\}} --- ссылка на номер
            \item \texttt{\textbackslash pageref\{метка\}} --- ссылка на страницу
        \end{itemize}
    \end{block}
    \vspace{0.5cm}
    
    \begin{exampleblock}{Соглашения об именовании}
        \begin{itemize}
            \item \texttt{fig:} --- для рисунков
            \item \texttt{tab:} --- для таблиц
            \item \texttt{eq:} --- для уравнений
            \item \texttt{sec:} --- для разделов
        \end{itemize}
    \end{exampleblock}
\end{frame}

% ==================== КОЛИЧЕСТВО ПРОГОНОВ ====================
\begin{frame}{Сколько прогонов нужно?}
    \begin{center}
        \begin{tikzpicture}[
            node distance=2cm,
            mystep/.style={rectangle, draw, fill=green!20, text width=3cm, align=center, minimum height=1.2cm}
        ]
            \node[mystep] (first) {\textbf{1-й прогон} \\ Создание меток \\ в .aux файле};
            \node[mystep, below=of first] (second) {\textbf{2-й прогон} \\ Чтение меток \\ и вставка ссылок};
            \node[mystep, below=of second] (third) {\textbf{3-й прогон} \\ (если изменилась \\ нумерация)};
            
            \draw[->, thick] (first) -- (second);
            \draw[->, thick] (second) -- (third);
        \end{tikzpicture}
    \end{center}
    
    \begin{alertblock}{Правило}
        Минимум \textbf{2 прогона} для корректных ссылок!
    \end{alertblock}
\end{frame}


% ==================== ЭКСПЕРИМЕНТ С LABEL ====================
\begin{frame}[fragile]{Эксперимент: label до caption}
    \begin{columns}[T]
        \begin{column}{0.48\textwidth}
            \begin{block}{Неправильно}
\begin{lstlisting}
\begin{figure}[h]
    \centering
    \label{fig:wrong}
    \caption{Risunok}
\end{figure}
\end{lstlisting}
                \vspace{0.2cm}
                {\color{red}\textbf{Результат:}} ссылка на номер \textbf{раздела}, а не рисунка!
            \end{block}
        \end{column}
        
        \begin{column}{0.48\textwidth}
            \begin{block}{Правильно}
\begin{lstlisting}
\begin{figure}[h]
    \centering
    \caption{Risunok}
    \label{fig:correct}
\end{figure}
\end{lstlisting}
                \vspace{0.2cm}
                {\color{green!50!black}\textbf{Результат:}} правильная ссылка на рисунок!
            \end{block}
        \end{column}
    \end{columns}
    
    \vspace{0.5cm}
    
    \begin{alertblock}{Запомните!}
        \texttt{\textbackslash label} всегда \textbf{после} \texttt{\textbackslash caption}!
    \end{alertblock}
\end{frame}

% ==================== ЭКСПЕРИМЕНТ С УРАВНЕНИЯМИ ====================
\begin{frame}[fragile]{Эксперимент: label для уравнений}
    \begin{columns}[T]
        \begin{column}{0.48\textwidth}
            \begin{block}{Правильно}
\begin{lstlisting}
\begin{equation}
    E = mc^2
    \label{eq:einstein}
\end{equation}
\end{lstlisting}
                {\color{green!50!black}\textbf{Работает!}}
            \end{block}
        \end{column}
        
        \begin{column}{0.48\textwidth}
            \begin{block}{Неправильно}
\begin{lstlisting}
\begin{equation}
    E = mc^2
\end{equation}
\label{eq:wrong}
\end{lstlisting}
                {\color{red}\textbf{Не работает!}}
            \end{block}
        \end{column}
    \end{columns}
    
    \vspace{0.5cm}
    
    \begin{alertblock}{Важно!}
        Для уравнений \texttt{\textbackslash label} должна быть \textbf{внутри} окружения \texttt{equation}!
    \end{alertblock}
    
    \begin{exampleblock}{Почему?}
        После \texttt{\textbackslash end\{equation\}} счётчик уже не активен
    \end{exampleblock}
\end{frame}

% ==================== РАЗДЕЛ 6: ПРАКТИЧЕСКИЕ ПРИМЕРЫ ====================
\section{Практические примеры}

\begin{frame}[fragile]{Полный пример документа}
    \begin{exampleblock}{Структура документа}
\begin{lstlisting}
\documentclass{article}
\usepackage[utf8]{inputenc}
\usepackage[russian]{babel}
\usepackage{graphicx}
\usepackage{lipsum}

\begin{document}
\section{Vvedenie}
\label{sec:intro}
\lipsum[1]

\begin{figure}[htbp]
    \centering
    \includegraphics[width=0.5\textwidth]{img.png}
    \caption{Primer izobrazheniya}
    \label{fig:example}
\end{figure}

Sm. ris.~\ref{fig:example} v razdele~\ref{sec:intro}.
\end{document}
\end{lstlisting}
    \end{exampleblock}
\end{frame}

% ==================== СОЗДАНИЕ СОБСТВЕННОГО ИЗОБРАЖЕНИЯ ====================
\begin{frame}{Создание собственного изображения в TikZ}
    \begin{center}
        \begin{tikzpicture}[scale=0.8]
            % Координатные оси
            \draw[->, thick] (-0.5,0) -- (6,0) node[right] {$x$};
            \draw[->, thick] (0,-0.5) -- (0,5) node[above] {$y$};
            
            % Сетка
            \draw[gray!30, very thin] (0,0) grid (5.5,4.5);
            
            % График функции y = x^2 / 4
            \draw[domain=0:4.2, smooth, variable=\x, blue, line width=2pt] 
                plot ({\x}, {0.25*\x*\x});
            
            % Точки на графике
            \fill[red] (0,0) circle (3pt) node[below left] {$O$};
            \fill[red] (2,1) circle (3pt) node[above right] {$A(2,1)$};
            \fill[red] (4,4) circle (3pt) node[above right] {$B(4,4)$};
            
            % Подписи осей
            \foreach \x in {1,2,3,4,5}
                \draw (\x,0.1) -- (\x,-0.1) node[below] {\tiny $\x$};
            \foreach \y in {1,2,3,4}
                \draw (0.1,\y) -- (-0.1,\y) node[left] {\tiny $\y$};
            
            % Формула функции
            \node[blue, above] at (3,3.5) {$y = \frac{x^2}{4}$};
        \end{tikzpicture}
    \end{center}
    
    \begin{block}{Описание}
        График параболы с отмеченными точками и координатной сеткой
    \end{block}
\end{frame}

% ==================== ПРИМЕР С TWOCOLUMN ====================
\begin{frame}[fragile]{Демонстрация twocolumn}
    \begin{exampleblock}{Код документа}
\begin{lstlisting}
\documentclass[twocolumn]{article}
\usepackage{graphicx}
\usepackage{lipsum}

\begin{document}
\section{Pervyj razdel}
\lipsum[1]

\begin{figure}[h]
    \centering
    \includegraphics[width=\linewidth]{image.png}
    \caption{Izobrazhenie v kolonke}
\end{figure}

\lipsum[2-3]
\end{document}
\end{lstlisting}
    \end{exampleblock}
\end{frame}

% ==================== ТАБЛИЦА СРАВНЕНИЯ ====================
\begin{frame}{Сравнительная таблица}
    \begin{table}
        \centering
        \caption{Сравнение параметров изображений}
        \label{tab:comparison}
        \small
        \begin{tabular}{@{}lccc@{}}
            \toprule
            \textbf{Параметр} & \textbf{Абсолютный} & \textbf{Относительный} & \textbf{Рекомендация} \\
            \midrule
            width & 5cm & 0.5\textbackslash textwidth & Относительный \\
            height & 3cm & 0.3\textbackslash textheight & Относительный \\
            scale & --- & 1.5 & Зависит \\
            angle & 45 & --- & Абсолютный \\
            \bottomrule
        \end{tabular}
    \end{table}
    
    \vspace{0.3cm}
    
    \begin{block}{Вывод}
        Предпочитайте относительные размеры для адаптивности!
    \end{block}
\end{frame}

% ==================== РАЗДЕЛ 7: ВЫВОДЫ ====================
\section{Выводы}

\begin{frame}{Основные выводы}
    \begin{enumerate}
        \item \textbf{Параметры изображений:}
        \begin{itemize}
            \item width, height, scale, angle работают независимо
            \item Можно комбинировать для нужного эффекта
        \end{itemize}
        
        \item \textbf{textwidth vs linewidth:}
        \begin{itemize}
            \item В одной колонке --- одинаковы
            \item В twocolumn --- linewidth меньше
        \end{itemize}
        
        \item \textbf{Плавающие объекты:}
        \begin{itemize}
            \item Спецификаторы [htbp] дают гибкость
            \item \LaTeX\ выбирает оптимальное размещение
        \end{itemize}
        
        \item \textbf{Метки и ссылки:}
        \begin{itemize}
            \item Нужно минимум 2 прогона компиляции
            \item label после caption для floats
            \item label внутри equation для уравнений
        \end{itemize}
    \end{enumerate}
\end{frame}

% ==================== РЕКОМЕНДАЦИИ ====================
\begin{frame}{Практические рекомендации}
    \begin{block}{Для изображений}
        \begin{itemize}
            \item Используйте \texttt{width=0.X\textbackslash linewidth}
            \item Добавляйте \texttt{keepaspectratio=true}
            \item Сохраняйте изображения в высоком разрешении
        \end{itemize}
    \end{block}
    
    \begin{block}{Для floats}
        \begin{itemize}
            \item Всегда используйте \texttt{\textbackslash centering}
            \item Добавляйте информативные подписи
            \item Используйте осмысленные метки (fig:name)
        \end{itemize}
    \end{block}
    
    \begin{block}{Для компиляции}
        \begin{itemize}
            \item Компилируйте 2-3 раза для ссылок
            \item Проверяйте предупреждения в логе
            \item Используйте latexmk для автоматизации
        \end{itemize}
    \end{block}
\end{frame}

% ==================== ЗАКЛЮЧЕНИЕ ====================
\begin{frame}{Заключение}
    \begin{block}{Что мы изучили}
        \begin{itemize}
            \item Параметры команды \texttt{\textbackslash includegraphics}
            \item Разницу между \texttt{textwidth} и \texttt{linewidth}
            \item Работу с плавающими объектами и спецификаторами
            \item Систему меток и ссылок в \LaTeX
            \item Важность порядка команд и количества прогонов
        \end{itemize}
    \end{block}
    
    \vspace{0.5cm}
    
    \begin{alertblock}{Ключевой вывод}
        \LaTeX\ требует понимания его логики работы, но взамен даёт профессиональное качество документов!
    \end{alertblock}
\end{frame}

% ==================== ФИНАЛЬНЫЙ СЛАЙД ====================
\begin{frame}
    \begin{center}
        {\Huge Спасибо за внимание!}
        
        \vspace{1.5cm}
        
    \end{center}
\end{frame}

\end{document}
